% =============================================================================
%  satSim — Physical and mathematical documentation of the ADCS simulator
%  Build:  pdflatex fisica_satSim.tex  (twice, for the table of contents and references)
% =============================================================================
\documentclass[11pt,a4paper]{article}

\usepackage[utf8]{inputenc}
\usepackage[T1]{fontenc}
\usepackage[english]{babel}
\usepackage{lmodern}
\usepackage{amsmath,amssymb,amsthm,bm}
\usepackage{geometry}
\usepackage{booktabs}
\usepackage{array}
\usepackage{xcolor}
\usepackage{tcolorbox}
\usepackage{enumitem}
\usepackage{listings}
\usepackage{hyperref}

\geometry{margin=2.5cm}
\hypersetup{colorlinks=true, linkcolor=blue!50!black, urlcolor=blue!50!black}
\setlength{\parskip}{0.4em}

% --- Environments -----------------------------------------------------------
\newtcolorbox{incode}[1][]{colback=gray!6, colframe=gray!50, fonttitle=\bfseries,
  title={In the code}, #1}
\newtcolorbox{note}[1][]{colback=blue!4, colframe=blue!40!black, fonttitle=\bfseries,
  title={#1}}
\newtcolorbox{improvement}{colback=green!5, colframe=green!40!black, fonttitle=\bfseries,
  title={Possible improvements}}
\newtheorem{prop}{Proposition}

\lstset{basicstyle=\ttfamily\small, columns=fullflexible, breaklines=true,
  frame=none, backgroundcolor=\color{gray!6}}

% --- Notation ---------------------------------------------------------------
\newcommand{\vect}[1]{\bm{#1}}
\newcommand{\w}{\vect{\omega}}
\newcommand{\q}{\vect{q}}
\newcommand{\h}{\vect{h}_{rw}}
\newcommand{\T}{\vect{T}}
\newcommand{\skewm}[1]{[\,#1\,]_\times}
\newcommand{\norm}[1]{\left\lVert #1 \right\rVert}
\newcommand{\code}[1]{\texttt{#1}}

\title{\textbf{satSim}\\[0.3em]
\large Physical and mathematical model of the attitude control simulator\\
for a 3U CubeSat with reaction wheels}
\author{Technical documentation --- version 0.1.0}
\date{\today}

\begin{document}
\maketitle

\begin{abstract}
This document describes in detail the physics and mathematics implemented
in the \emph{satSim} simulator: the rigid-body model of the satellite, the
dynamics of the reaction wheels and their power consumption, the Euler
equations with gyroscopic coupling, quaternion kinematics, the models of the
environmental disturbance torques (gravity gradient, solar radiation
pressure, magnetic torque), the fourth-order Runge--Kutta numerical
integrator and the quaternion PD controller used as the baseline.
For each model the derivation, the assumptions, the orders of magnitude,
the correspondence with the code and the directions for improvement are given.
The final goal of the project is to build on the engine a \emph{Gymnasium}
environment to train a Reinforcement Learning agent (PPO/SAC).
\end{abstract}

\tableofcontents
\newpage

% =============================================================================
\section{Introduction and working principle}
% =============================================================================

\subsection{The attitude control problem}
The \emph{attitude} of a satellite is its orientation in space relative to
a reference frame. The ADCS subsystem (\emph{Attitude Determination
and Control System}) has to measure the attitude and bring it to / keep it at a
desired value (\emph{target}), for example to point an antenna, a camera or the
solar panels.

In space there is nothing to ``push against'': to rotate the satellite one
must either expel mass (thrusters), interact with an external field
(magnetorquers), or \textbf{exchange angular momentum} with internal moving
parts. \emph{Reaction wheels} belong to the last category.

\subsection{Physical principle of reaction wheels}
A reaction wheel is a flywheel driven by a brushless electric motor fixed to
the structure. If the motor applies a torque $T_{rw}$ to the rotor, by Newton's
third law the rotor applies the torque $-T_{rw}$ to the structure. The whole
system (body + wheels), in the absence of external torques, conserves its total
angular momentum:
\begin{equation}
  \vect{H} = \underbrace{J\,\w}_{\text{body}} + \underbrace{\h}_{\text{wheels}},
  \qquad
  \left.\frac{d\vect{H}}{dt}\right|_{\mathcal{I}} = \T_{ext}.
  \label{eq:Htot}
\end{equation}
Spinning up a wheel in one direction therefore rotates the body in the opposite
direction: angular momentum is \emph{transferred}, not created. Two fundamental
consequences follow, which the RL agent has to learn to manage:
\begin{enumerate}[nosep]
  \item the wheels can only \emph{redistribute} angular momentum; external
        disturbance torques instead \emph{add} it to the system;
  \item the accumulated momentum makes the wheel speeds grow up to the
        physical limit (\textbf{saturation}); from then on control of the
        corresponding axis is lost until the wheel is ``unloaded'' with an
        external actuator (\emph{momentum dumping}).
\end{enumerate}

\subsection{The system of equations solved}
The simulator numerically integrates a system of nonlinear ordinary
differential equations
\begin{equation}
  \dot{\vect{x}} = f(t,\vect{x},\vect{u}),\qquad
  \vect{x} = \big[\,\q^T,\ \w^T,\ \vect{\Omega}^T,\ E\,\big]^T \in \mathbb{R}^{7+N+1},
  \label{eq:state}
\end{equation}
where $\q\in\mathbb{R}^4$ is the attitude quaternion, $\w\in\mathbb{R}^3$ the
angular velocity of the body, $\vect\Omega\in\mathbb{R}^N$ the speeds of the $N$
wheels, $E$ the electrical energy consumed and $\vect{u}\in\mathbb{R}^N$ the
vector of motor torques commanded to the wheels. The five ``equations'' of the
model are:
\begin{center}
\begin{tabular}{cll}
\toprule
\# & Content & Section\\
\midrule
1 & Inertia matrix $J$ of the 3U CubeSat & \S\ref{sec:inertia}\\
2 & Wheel dynamics $\dot\Omega = T_{rw}/I_{rw}$, saturation, power & \S\ref{sec:wheels}\\
3 & Euler equations with gyroscopic coupling & \S\ref{sec:euler}\\
4 & Quaternion kinematics & \S\ref{sec:quat}\\
5 & Environmental disturbance torques & \S\ref{sec:disturbances}\\
\bottomrule
\end{tabular}
\end{center}

% =============================================================================
\section{Reference frames and notation}\label{sec:frames}
% =============================================================================

\subsection{Reference frames}
\begin{itemize}
  \item \textbf{Inertial} $\mathcal{I}$: geocentric ECI frame (\emph{Earth
        Centered Inertial}), with $z_I$ along the Earth's polar axis. In the
        simulator it also coincides with the \textbf{target attitude} (identity
        quaternion): the controller tries to align the satellite axes with
        those of $\mathcal{I}$.
  \item \textbf{Body} $\mathcal{B}$: attached to the satellite, origin at the
        centre of mass, axes parallel to the edges of the box; $z_B$ is the
        long axis (30 cm).
\end{itemize}

\subsection{Notation}
\begin{center}
\begin{tabular}{>{$}l<{$}l}
\toprule
\text{Symbol} & Meaning\\
\midrule
\q = [q_0, q_1, q_2, q_3]^T & unit quaternion, scalar part $q_0$ first\\
\q_v = [q_1,q_2,q_3]^T & vector part of the quaternion\\
R(\q) & rotation matrix $\mathcal{B}\to\mathcal{I}$: $\vect{v}_I = R(\q)\,\vect{v}_B$\\
\w & angular velocity of $\mathcal{B}$ relative to $\mathcal{I}$, expressed in $\mathcal{B}$ [rad/s]\\
J & inertia tensor of the satellite in $\mathcal{B}$ [kg\,m$^2$]\\
\Omega_i & axial speed of wheel $i$ [rad/s]\\
A\in\mathbb{R}^{3\times N} & distribution matrix: column $i$ = axis of wheel $i$ in $\mathcal{B}$\\
\skewm{\vect{a}} & skew-symmetric matrix such that $\skewm{\vect a}\vect b = \vect a\times\vect b$\\
\bottomrule
\end{tabular}
\end{center}
All quantities are in SI units, except the wheel speeds shown in RPM in the
GUI ($1\ \text{RPM} = 2\pi/60$ rad/s).

% =============================================================================
\section{Equation 1 --- Inertia matrix}\label{sec:inertia}
% =============================================================================

\subsection{Definition}
The inertia tensor of a rigid body about its centre of mass is
\begin{equation}
  J = \int_V \left(\norm{\vect\rho}^2 I_3 - \vect\rho\,\vect\rho^T\right) dm,
\end{equation}
where $\vect\rho$ is the position of the mass element $dm$. The diagonal
elements are the \emph{moments of inertia}, the off-diagonal ones the
\emph{products of inertia} (with a negative sign: $J_{xy} = -\int xy\,dm$). $J$ is
symmetric and positive definite; its eigenvectors are the \emph{principal
axes of inertia}.

\subsection{3U CubeSat as a homogeneous box}
For a homogeneous box of mass $m$ and sides $a,b,c$ along $x,y,z$,
the integral gives
\begin{equation}
  J_{xx}=\frac{m}{12}(b^2+c^2),\qquad
  J_{yy}=\frac{m}{12}(a^2+c^2),\qquad
  J_{zz}=\frac{m}{12}(a^2+b^2).
\end{equation}
With $m=3$ kg and $a\times b\times c = 0.1\times0.1\times0.3$ m:
$J_{xx}=J_{yy}=0.025$ kg\,m$^2$, $J_{zz}=0.005$ kg\,m$^2$. The satellite is
\emph{prolate}: the $z_B$ axis has the minimum inertia (five times ``easier''
to rotate than the others).

\subsection{Products of inertia}
A real satellite is not homogeneous (batteries, boards, payload), so the
geometric axes do not exactly coincide with the principal ones. Small
products of inertia are introduced:
\begin{equation}
J=\begin{bmatrix}
0.025 & 2\cdot10^{-5} & -1\cdot10^{-5}\\
2\cdot10^{-5} & 0.025 & 1.5\cdot10^{-5}\\
-1\cdot10^{-5} & 1.5\cdot10^{-5} & 0.005
\end{bmatrix}\ \text{kg\,m}^2 .
\end{equation}
They create coupling between the axes: a torque about $x$ also produces a
(very small) acceleration about $y$ and $z$. The code checks that all
eigenvalues are positive.

\begin{incode}
\code{physics\_engine.cubesat\_inertia(mass, size, products)}. $J^{-1}$ is
computed only once in the constructor of \code{SatelliteEngine}, because $J$
is constant.
\end{incode}

\begin{improvement}
\begin{itemize}[nosep]
  \item Build $J$ as the sum of components (structure, batteries, wheels,
        payload) with the parallel axis (Huygens--Steiner) theorem:
        $J = \sum_k \left[J_k + m_k(\norm{\vect d_k}^2 I_3 - \vect d_k\vect d_k^T)\right]$.
  \item Explicitly subtract the axial inertia of the wheels,
        $J_{body} = J_{tot} - I_{rw}\sum_i \vect a_i\vect a_i^T$ (see \S\ref{sec:wheels-note}).
  \item Randomise $J$ ($\pm 10\%$) during RL training
        (\emph{domain randomization}) to obtain a robust policy.
  \item Deployable solar panels: a much larger $J$, no longer
        ``prolate''; possibly flexible appendages.
\end{itemize}
\end{improvement}

% =============================================================================
\section{Equation 2 --- Reaction wheels}\label{sec:wheels}
% =============================================================================

\subsection{Axial dynamics}
Each wheel is an axisymmetric rotor with axial moment of inertia
$I_{rw}$. The angular momentum balance along the spin axis gives
\begin{equation}
  I_{rw}\,\dot\Omega_i = T_{rw,i}
  \qquad\Longleftrightarrow\qquad
  \frac{d\Omega_i}{dt} = \frac{T_{rw,i}}{I_{rw}}.
  \label{eq:rw}
\end{equation}
The angular momentum of the array, expressed in $\mathcal{B}$, and the reaction
torque on the body are
\begin{equation}
  \h = A\,I_{rw}\,\vect\Omega,
  \qquad
  \T_{rw,action} = -A\,\vect{T}_{rw}.
  \label{eq:hrw}
\end{equation}

\subsection{Note on the definition of $\Omega$}\label{sec:wheels-note}
In equation~\eqref{eq:rw} $\Omega_i$ must be interpreted as the \emph{absolute}
(inertial) speed of the rotor about its own axis. With this choice
\eqref{eq:rw} is exact for a symmetric rotor and the total angular momentum
\eqref{eq:Htot} is conserved exactly (see
Proposition~\ref{prop:H}). The speed measured by an encoder would be the
\emph{relative} one, $\Omega_i^{rel} = \Omega_i - \vect a_i^T\w$; the difference
($\sim10^{-2}$ rad/s against $\sim10^2$ rad/s) is negligible for
visualisation but should be considered when modelling real sensors.

\subsection{Geometric configurations}
\begin{itemize}
  \item \textbf{Orthogonal} (3 wheels): $A = I_3$. Each wheel controls one axis.
  \item \textbf{Pyramid} (4 wheels, default): the axes are tilted by
        $\beta$ from $z_B$ and distributed in azimuth
        $\varphi_i\in\{45^\circ,135^\circ,225^\circ,315^\circ\}$:
        \begin{equation}
          \vect a_i = [\sin\beta\cos\varphi_i,\ \sin\beta\sin\varphi_i,\ \cos\beta]^T,
          \qquad \beta = 54.74^\circ = \arccos(1/\sqrt3).
        \end{equation}
        With this $\beta$ the torque/momentum capacity is the same on the three
        axes. The configuration is \emph{redundant}: with one failed wheel the
        other three still have rank 3.
\end{itemize}

\subsection{Torque allocation}
The controller produces a desired body torque $\T_{cmd}\in\mathbb{R}^3$.
One has to find $\vect T_{rw}\in\mathbb{R}^N$ such that $-A\vect T_{rw}=\T_{cmd}$.
For $N=4$ the system is underdetermined; the minimum-norm solution is
chosen:
\begin{equation}
  \vect T_{rw} = -A^{+}\,\T_{cmd},\qquad A^{+} = A^T(AA^T)^{-1},
\end{equation}
where $A^+$ is the Moore--Penrose pseudo-inverse. Any solution can be
written as $-A^+\T_{cmd} + \vect n$, with $\vect n\in\ker A$ (the \emph{null
space}, of dimension $N-3=1$): components along $\vect n$ change the wheel
speeds \emph{without} producing torque on the body.

\subsection{Saturations}
Two physical limits:
\begin{align}
  |T_{rw,i}| &\le T_{max} = 2\ \text{mN\,m} &&\text{(maximum motor current)},\\
  |\Omega_i| &\le \Omega_{max} = 6000\ \text{RPM}\approx 628\ \text{rad/s} &&\text{(mechanical/back-EMF limit)}.
\end{align}
Since the torque is held constant over the whole integration step
$\Delta t$ (\S\ref{sec:rk4}), equation~\eqref{eq:rw} has the exact solution
$\Omega_i(t+\Delta t) = \Omega_i(t) + T_{rw,i}\Delta t/I_{rw}$. The command is
therefore limited \emph{before} integration:
\begin{equation}
  \frac{(-\Omega_{max}-\Omega_i)\,I_{rw}}{\Delta t}
  \ \le\ T_{rw,i}\ \le\
  \frac{(\Omega_{max}-\Omega_i)\,I_{rw}}{\Delta t}.
\end{equation}
This way the wheel reaches exactly $\Omega_{max}$ without ever exceeding it,
and there is no need to ``cut'' the state afterwards
(\emph{clipping}), an operation that would destroy angular momentum in a
non-physical way. At saturation the wheel accepts torque only in the direction
that slows it down.

The maximum angular momentum a wheel can store is
$h_{max} = I_{rw}\Omega_{max} = 1.5\cdot10^{-5}\cdot628 \approx 9.4$ mN\,m\,s.

\subsection{Power consumption model}
The power drawn by wheel $i$ is modelled as
\begin{equation}
  P_i = k_1\,|T_{rw,i}| + k_2\,|T_{rw,i}\,\Omega_i| + P_{static}.
  \label{eq:power}
\end{equation}
\begin{itemize}
  \item $k_1|T|$: in a brushless motor $T = k_t\,i$; the losses (resistive,
        in the driver and in the windings) grow with the current and hence with
        the torque. The real dependence is quadratic ($R i^2$); here it is
        linearised.
  \item $k_2|T\Omega|$: $T\Omega$ is the mechanical power transferred to the rotor;
        $k_2 = 1/\eta \approx 1.2$ accounts for the efficiency. The absolute
        value implies that braking does \emph{not} regenerate energy.
  \item $P_{static}$: control electronics, encoder, bearing losses.
\end{itemize}
The energy $E = \int_0^t \sum_i P_i\,d\tau$ is part of the state vector and is
integrated with RK4 together with the dynamics: $\dot E = \sum_i P_i$. This keeps
it accurate even when $\Omega_i$ changes a lot within a step.

\begin{incode}
Class \code{ReactionWheelArray}: \code{\_spin\_axes()} (matrix $A$),
\code{limit\_torque()} (saturations), \code{power()} (eq.~\ref{eq:power}),
\code{momentum\_body()} (eq.~\ref{eq:hrw}); allocation in
\code{SatelliteEngine.allocate()}.
\end{incode}

\begin{improvement}
\begin{itemize}[nosep]
  \item \textbf{Friction}: $T_f = -c_v\Omega - T_c\,\mathrm{sign}(\Omega)$
        (viscous + Coulomb), with the typical ``dead zone'' around
        $\Omega=0$ (\emph{zero-crossing}), a source of real disturbances.
  \item \textbf{Motor dynamics}: first-order lag
        $\tau_m\dot T = T_{cmd}-T$ and command quantisation.
  \item \textbf{Speed-dependent torque limit} (back-EMF):
        $T_{max}(\Omega) = T_{max}\,(1-|\Omega|/\Omega_{no\text{-}load})$.
  \item \textbf{More realistic power}: $P = R\,(T/k_t)^2 + T\Omega/\eta + P_s$,
        possibly with regenerative braking.
  \item Use the \textbf{null space} (4 wheels) to keep the speeds away from
        zero and from saturation.
  \item Static/dynamic rotor imbalance (\emph{micro-vibrations}).
\end{itemize}
\end{improvement}

% =============================================================================
\section{Equation 3 --- Euler equations with gyroscopic coupling}\label{sec:euler}
% =============================================================================

\subsection{Derivation}
Start from the angular momentum balance in an inertial frame:
\begin{equation}
  \left.\frac{d\vect H}{dt}\right|_{\mathcal I} = \T_{ext}.
\end{equation}
For a vector expressed in a frame rotating with angular velocity $\w$ the
\emph{transport theorem} holds:
\begin{equation}
  \left.\frac{d\vect H}{dt}\right|_{\mathcal I} =
  \left.\frac{d\vect H}{dt}\right|_{\mathcal B} + \w\times\vect H .
\end{equation}
Substituting $\vect H = J\w + \h$, with $J$ constant in $\mathcal{B}$:
\begin{equation}
  J\dot\w + \dot{\vect h}_{rw} + \w\times(J\w+\h) = \T_{dist}.
\end{equation}
From \eqref{eq:rw}--\eqref{eq:hrw}, $\dot{\vect h}_{rw} = A I_{rw}\dot{\vect\Omega} = A\vect T_{rw}$.
Moving this term to the right-hand side gives the implemented equation:
\begin{equation}
  \boxed{\;J\,\dot\w = \T_{tot} - \w\times\left(J\w + \h\right),
  \qquad \T_{tot} = -A\vect T_{rw} + \T_{dist}\;}
  \label{eq:euler}
\end{equation}
and therefore $\dot\w = J^{-1}\left[\T_{tot} - \w\times(J\w+\h)\right]$.

\subsection{Physical interpretation of the terms}
\begin{itemize}
  \item $\w\times J\w$: nonlinear coupling of the rigid body.
        It is responsible, for example, for free precession (Euler--Poinsot
        motion) and for the instability of rotation about the axis of
        intermediate inertia (\emph{Dzhanibekov effect}).
  \item $\w\times\h$: \emph{gyroscopic stiffness} of the wheels. If the body
        rotates while the wheels store momentum, a torque appears about an
        axis perpendicular to both. It becomes important when the wheels
        spin fast.
\end{itemize}

\subsection{Conservation of angular momentum}
\begin{prop}\label{prop:H}
If $\T_{dist}=\vect 0$, the total angular momentum expressed in the inertial
frame, $\vect H_I = R(\q)\,(J\w + \h)$, is constant for any command
$\vect T_{rw}(t)$.
\end{prop}
\begin{proof}
Let $\vect H_B = J\w+\h$. From \eqref{eq:euler} and $\dot{\vect h}_{rw}=A\vect T_{rw}$:
$\dot{\vect H}_B = -A\vect T_{rw} - \w\times\vect H_B + A\vect T_{rw} = -\w\times\vect H_B$.
Moreover $\dot R = R\,\skewm{\w}$. Hence
$\dot{\vect H}_I = R\skewm{\w}\vect H_B + R\dot{\vect H}_B = R(\w\times\vect H_B - \w\times\vect H_B) = \vect 0$.
\end{proof}
This property is checked numerically in the tests (error $<10^{-9}$ N\,m\,s
after 100 s with random torques) and is the most effective check of the sign
consistency between wheels, Euler and allocation.

Similarly, without wheels or torques, the kinetic energy
$K = \tfrac12\w^TJ\w$ is conserved, since $\dot K = \w^TJ\dot\w = -\w^T(\w\times J\w) = 0$.

\begin{incode}
\code{SatelliteEngine.\_derivatives()}: \code{w\_dot = J\_inv @ (T\_total - cross3(w, J @ w + h))}.
\end{incode}

% =============================================================================
\section{Equation 4 --- Quaternion kinematics}\label{sec:quat}
% =============================================================================

\subsection{Why quaternions}
Euler angles (3 parameters) have singularities
(\emph{gimbal lock}): for some attitudes the map from angular velocities to
angle derivatives diverges. No 3-parameter parametrisation of
$SO(3)$ is globally singularity-free. Unit quaternions use 4
parameters with the constraint $\norm{\q}=1$ and are singularity-free, as well
as computationally cheap (no trigonometric functions).

\subsection{Definition and algebra}
A unit quaternion representing a rotation by an angle $\theta$
about the unit axis $\vect e$ is
\begin{equation}
  \q = \begin{bmatrix}\cos(\theta/2)\\ \sin(\theta/2)\,\vect e\end{bmatrix}.
\end{equation}
The Hamilton product is
\begin{equation}
  \vect p\otimes\q =
  \begin{bmatrix}
    p_0q_0 - \vect p_v\cdot\q_v\\
    p_0\q_v + q_0\vect p_v + \vect p_v\times\q_v
  \end{bmatrix},
\end{equation}
and the conjugate $\q^* = [q_0, -\q_v^T]^T$ represents the inverse rotation.
The rotation of a vector is written $[0,\vect v_I] = \q\otimes[0,\vect v_B]\otimes\q^*$,
equivalent to $\vect v_I = R(\q)\vect v_B$ with
\begin{equation}
R(\q)=\begin{bmatrix}
1-2(q_2^2+q_3^2) & 2(q_1q_2-q_0q_3) & 2(q_1q_3+q_0q_2)\\
2(q_1q_2+q_0q_3) & 1-2(q_1^2+q_3^2) & 2(q_2q_3-q_0q_1)\\
2(q_1q_3-q_0q_2) & 2(q_2q_3+q_0q_1) & 1-2(q_1^2+q_2^2)
\end{bmatrix}.
\end{equation}

\subsection{Kinematic equation}
With $\w$ expressed in body axes, the quaternion derivative is
\begin{equation}
  \dot\q = \frac12\,\q\otimes\begin{bmatrix}0\\ \w\end{bmatrix}
  = \frac12\begin{bmatrix}
    -q_1\omega_x - q_2\omega_y - q_3\omega_z\\
    \ \ q_0\omega_x - q_3\omega_y + q_2\omega_z\\
    \ \ q_3\omega_x + q_0\omega_y - q_1\omega_z\\
    -q_2\omega_x + q_1\omega_y + q_0\omega_z
  \end{bmatrix}.
  \label{eq:qdot}
\end{equation}
\emph{Idea of the derivation}: in an interval $dt$ the body performs an
infinitesimal rotation $d\q \approx [1, \tfrac12\w\,dt]$ about its own (body)
axes, which is composed \emph{on the right}: $\q(t+dt) = \q(t)\otimes d\q$. Hence
$\dot\q = \tfrac12\q\otimes[0,\w]$.

\subsection{Normalisation}
Equation \eqref{eq:qdot} preserves the norm exactly
($\tfrac{d}{dt}\norm{\q}^2 = 2\q\cdot\dot\q = 0$), but the numerical integrator
does not: the truncation error makes $\norm{\q}$ drift slowly. After every
step the quaternion is therefore projected back onto the unit sphere:
$\q\leftarrow\q/\norm{\q}$.

\subsection{Attitude error}
Given the target $\q_t$, the error is the rotation that takes the target to the
current attitude:
\begin{equation}
  \q_e = \q_t^*\otimes\q,\qquad
  \theta_{err} = 2\arccos|q_{e,0}|.
\end{equation}
Since $\q$ and $-\q$ represent the same attitude (double cover of
$SO(3)$), the sign with $q_{e,0}\ge0$ is chosen: thus $\theta_{err}\le180^\circ$
and the controller always rotates the shortest way, avoiding \emph{unwinding}
(an almost complete, useless turn).

\begin{incode}
\code{quaternion.py}: \code{quat\_mult}, \code{quat\_to\_dcm}, \code{quat\_error},
\code{quat\_angle\_deg}. Normalisation in \code{SatelliteEngine.step()}.
\end{incode}

% =============================================================================
\section{Equation 5 --- Environmental disturbance torques}\label{sec:disturbances}
% =============================================================================

The total torque acting on the body is
\begin{equation}
  \T_{tot} = \T_{rw,action} + \T_{gg} + \T_{srp} + \T_{mag}\ (+\ \T_{manual}).
\end{equation}
The disturbances depend on the orbital position and on the attitude, and are
recomputed at every stage of the RK4 integrator.

\subsection{Orbit model}
Circular Keplerian orbit of radius $r = R_\oplus + 500$ km, inclination
$i = 51.6^\circ$, right ascension of the ascending node $\Omega_{RAAN}$:
\begin{equation}
  \vect r_I(t) = R_3(-\Omega_{RAAN})\,R_1(-i)\;r\begin{bmatrix}\cos u\\ \sin u\\ 0\end{bmatrix},
  \qquad u = u_0 + n t,\quad n=\sqrt{\mu/r^3},
\end{equation}
with $\mu = 3.986\cdot10^{14}$ m$^3$/s$^2$. The period is
$T_{orb} = 2\pi/n \approx 94.5$ min.

\subsection{Gravity gradient}
\paragraph{Physical origin.} The gravitational field decreases as $1/r^2$: the
part of the satellite closest to the Earth is pulled slightly more than the
far part. If the body is not symmetric about the local vertical, this
difference produces a torque.

\paragraph{Derivation.} The torque is $\T = \int \vect\rho\times d\vect F$ with
$d\vect F = -\mu\,(\vect r+\vect\rho)/\norm{\vect r+\vect\rho}^3\,dm$.
Expanding to first order in $\rho/r$ and using
$\int\vect\rho\,dm = 0$ (origin at the centre of mass):
\begin{equation}
  \T_{gg} = \frac{3\mu}{r^3}\;\hat{\vect r}_B\times\left(J\,\hat{\vect r}_B\right),
  \qquad
  \hat{\vect r}_B = R(\q)^T\,\frac{\vect r_I}{r}.
\end{equation}

\paragraph{Properties.}
\begin{itemize}[nosep]
  \item $\T_{gg}=\vect 0$ if $\hat{\vect r}_B$ is an eigenvector of $J$ (principal
        axis aligned with the vertical).
  \item It tends to align the axis of \emph{minimum} inertia with the vertical:
        the principle of passive gravity-gradient stabilisation.
  \item Order of magnitude: $3\mu/r^3\approx 3.4\cdot10^{-6}$ s$^{-2}$,
        $\Delta J\approx 0.02$ kg\,m$^2$ $\Rightarrow$ $|\T_{gg}|\lesssim 3.5\cdot10^{-8}$ N\,m.
\end{itemize}

\subsection{Solar radiation pressure (SRP)}
\paragraph{Physical origin.} Photons carry momentum $p=E/c$.
At 1 AU the solar flux $\Phi\approx 1361$ W/m$^2$ exerts a pressure
$P_\odot = \Phi/c \approx 4.56\cdot10^{-6}$ N/m$^2$.

\paragraph{Flat-face model.} The satellite is described by its 6 faces
(normal $\hat{\vect n}_k$, area $A_k$, centre $\vect r_k$). With
$\hat{\vect s}$ the Sun direction in body axes and
$\cos\theta_k = \hat{\vect n}_k\cdot\hat{\vect s}$, every lit face
($\cos\theta_k>0$) receives the force
\begin{equation}
  \vect F_k = -P_\odot A_k\cos\theta_k
  \left[(1-\rho_s)\,\hat{\vect s} + 2\left(\rho_s\cos\theta_k + \frac{\rho_d}{3}\right)\hat{\vect n}_k\right],
\end{equation}
where $\rho_s,\rho_d$ are the fractions of light reflected specularly and
diffusely (the fraction $1-\rho_s-\rho_d$ is absorbed). The three contributions
are: absorption (push along $\hat{\vect s}$), specular reflection
(push along the normal, factor 2 for the momentum reversal),
Lambertian diffuse reflection (factor $2/3$ along the normal). The torque is
\begin{equation}
  \T_{srp} = \sum_{k:\cos\theta_k>0} (\vect r_k - \vect r_{cm})\times\vect F_k .
\end{equation}
The torque arises from the \textbf{offset between the centre of pressure and the centre of mass}:
if the centre of mass coincided with the geometric centre, the torques would
cancel by symmetry. In the simulator $\vect r_{cm} = [2,-1,10]$ mm.

\paragraph{Eclipse.} Cylindrical shadow model: the satellite is in shadow if
$\vect r\cdot\hat{\vect s}<0$ and its distance from the Earth--Sun axis is less
than $R_\oplus$:
\begin{equation}
  \text{eclipse}\iff \vect r\cdot\hat{\vect s} < 0\ \wedge\
  \norm{\vect r - (\vect r\cdot\hat{\vect s})\hat{\vect s}} < R_\oplus .
\end{equation}
In eclipse $\T_{srp}=\vect 0$. Order of magnitude:
$P_\odot\cdot0.03\ \text{m}^2\cdot0.01\ \text{m}\approx 10^{-9}$ N\,m.

\subsection{Residual magnetic torque}
\paragraph{Physical origin.} Currents in the circuits and magnetised materials
give the satellite a residual magnetic dipole moment $\vect m$
(typically $10^{-3}$--$10^{-2}$ A\,m$^2$ for a CubeSat). In a field
$\vect B$ a dipole experiences the torque $\vect m\times\vect B$ (like a
compass needle).

\paragraph{Field model.} Dipole Earth field, with dipole axis
$\hat{\vect m}_\oplus = -\hat{\vect z}_I$:
\begin{equation}
  \vect B_I(\vect r) = B_0\left(\frac{R_\oplus}{r}\right)^3
  \left[3(\hat{\vect m}_\oplus\cdot\hat{\vect r})\hat{\vect r} - \hat{\vect m}_\oplus\right],
  \qquad B_0 = 3.12\cdot10^{-5}\ \text{T},
\end{equation}
\begin{equation}
  \T_{mag} = \vect m\times\left(R(\q)^T\vect B_I\right).
\end{equation}
At 500 km $|\vect B|\approx 25$--$50\ \mu$T, hence
$|\T_{mag}|\sim 10^{-7}$ N\,m: it is typically the \textbf{dominant disturbance}
for a CubeSat in low Earth orbit.

\subsection{Manual torque}
Constant body torque set from the GUI ($\pm5$ mN\,m) or a 1 s impulse.
It is deliberately several orders of magnitude larger than the natural
disturbances, to make the controller response and the wheel saturation visible.

\subsection{Comparison of orders of magnitude}
\begin{center}
\begin{tabular}{lcc}
\toprule
Torque & Order of magnitude [N\,m] & Momentum accumulated in one orbit [N\,m\,s]\\
\midrule
Gravity gradient & $10^{-8}$ & $\lesssim 10^{-4}$ (partly periodic)\\
Solar pressure & $10^{-9}$ & $\sim 10^{-6}$\\
Residual magnetic & $10^{-7}$ & $\sim 10^{-4}$--$10^{-3}$\\
Wheel capacity (action) & $2\cdot10^{-3}$ & $h_{max}\approx 9.4\cdot10^{-3}$ per wheel\\
\bottomrule
\end{tabular}
\end{center}
The \emph{secular} (non-periodic) components of the disturbances accumulate in
the wheels orbit after orbit: without \emph{momentum dumping} the wheels
sooner or later saturate.

\begin{incode}
\code{environment.OrbitalEnvironment}: \code{position()}, \code{in\_eclipse()},
\code{magnetic\_field\_inertial()}, \code{srp\_torque()},
\code{disturbance\_torques()}.
\end{incode}

\begin{improvement}
\begin{itemize}[nosep]
  \item \textbf{Aerodynamic torque}: at 500 km it is of the same order as SRP and
        gravity gradient: $\vect F_k = -\tfrac12\rho_{atm} v^2 C_D A_k \max(0,\hat{\vect n}_k\cdot\hat{\vect v})\,\hat{\vect v}$,
        with density from an exponential model or NRLMSISE-00.
  \item \textbf{IGRF magnetic field} (tilted dipole + higher harmonics,
        Earth rotation); residual dipole varying with the current drawn.
  \item \textbf{Orbit}: propagation with $J_2$ perturbation or SGP4 from real
        TLEs; apparent motion of the Sun (ephemerides); conical shadow
        (penumbra).
  \item LVLH orbital frame and nadir-pointing target (more realistic
        for Earth observation).
  \item Dumping actuators: \textbf{magnetorquers} with the law
        $\vect m_{cmd} = \frac{k}{\norm{\vect B}^2}\,\vect B\times\Delta\vect h$.
\end{itemize}
\end{improvement}

% =============================================================================
\section{Numerical integration}\label{sec:rk4}
% =============================================================================

\subsection{Fourth-order Runge--Kutta}
The system $\dot{\vect x} = f(t,\vect x,\vect u)$ is integrated with fixed-step
RK4, $\Delta t = 0.05$ s:
\begin{align}
  \vect k_1 &= f(t_n,\ \vect x_n,\ \vect u_n),\\
  \vect k_2 &= f(t_n+\tfrac{\Delta t}{2},\ \vect x_n+\tfrac{\Delta t}{2}\vect k_1,\ \vect u_n),\\
  \vect k_3 &= f(t_n+\tfrac{\Delta t}{2},\ \vect x_n+\tfrac{\Delta t}{2}\vect k_2,\ \vect u_n),\\
  \vect k_4 &= f(t_n+\Delta t,\ \vect x_n+\Delta t\,\vect k_3,\ \vect u_n),\\
  \vect x_{n+1} &= \vect x_n + \frac{\Delta t}{6}\left(\vect k_1+2\vect k_2+2\vect k_3+\vect k_4\right).
\end{align}
The local truncation error is $O(\Delta t^5)$, the global one $O(\Delta t^4)$.
After the step: quaternion normalisation.

\subsection{Design choices}
\begin{itemize}
  \item \textbf{Zero-order hold}: the command $\vect u_n$ stays constant for the
        whole step. It reproduces an on-board computer that updates the
        actuators at a fixed rate (here 20 Hz) and is exactly the semantics of
        \code{env.step(action)} in Gymnasium.
  \item \textbf{Fixed step}: deterministic and reproducible simulation, constant
        cost per step (important for RL).
  \item \textbf{Choice of $\Delta t$}: the time scales of the system (controller
        period $2\pi/\omega_n\approx 16$ s, orbital period 94 min)
        are much larger than $\Delta t$. The tightest constraint is the
        tumbling rate: for $|\w|\Delta t \ll 1$ with
        $|\w|\lesssim 1$ rad/s there is ample margin.
  \item The disturbances are evaluated at every stage $\vect k_j$, with the
        intermediate attitude and time.
\end{itemize}

\begin{improvement}
\begin{itemize}[nosep]
  \item \textbf{Geometric integrators} (e.g. Crouch--Grossman, Lie methods
        or the exponential update $\q_{n+1} = \q_n\otimes\exp(\tfrac12\w\Delta t)$)
        that preserve $\norm{\q}=1$ exactly.
  \item Physics sub-steps smaller than the control step (e.g. physics at
        100 Hz, control at 10 Hz), also useful for RL (\emph{action repeat}).
  \item Performance: vectorise over several parallel environments or compile
        with \code{numba} (today $\approx1.5$ ms/step in NumPy).
\end{itemize}
\end{improvement}

% =============================================================================
\section{Quaternion PD controller}\label{sec:pd}
% =============================================================================

\subsection{Control law}
As the classical baseline for the RL agent, the satellite is stabilised by the law
\begin{equation}
  \T_{cmd} = -J\left(K_p\,\q_{e,v} + K_d\,\w\right) + \w\times(J\w+\h),
  \label{eq:pd}
\end{equation}
followed by the allocation $\vect T_{rw} = -A^+\T_{cmd}$ and by the saturations.
\begin{itemize}
  \item The first term is an \emph{inertia-normalised} PD: multiplying
        by $J$ gives every axis the same closed-loop dynamics.
  \item The second term \emph{compensates} the gyroscopic coupling
        (partial \emph{feedback linearisation}).
\end{itemize}

\subsection{Choice of the gains}
With the compensation, and without disturbances or saturations, substituting
\eqref{eq:pd} into \eqref{eq:euler} gives
$\dot\w = -K_p\q_{e,v} - K_d\w$. For small angles $\q_{e,v}\approx\vect\theta/2$
and $\w\approx\dot{\vect\theta}$, so every axis is a second-order
oscillator:
\begin{equation}
  \ddot\theta + K_d\dot\theta + \frac{K_p}{2}\theta = 0
  \quad\Longrightarrow\quad
  K_p = 2\omega_n^2,\qquad K_d = 2\zeta\omega_n .
\end{equation}
Default values: $\omega_n = 0.4$ rad/s, $\zeta = 0.9$ (almost critically
damped), giving a 2\% settling time of about
$4/(\zeta\omega_n)\approx 11$ s in the linear regime. For large initial
errors the requested torque exceeds $T_{max}$ and the response is slower.

\subsection{Stability (Lyapunov)}
Consider
\begin{equation}
  V = 2K_p\,(1-q_{e,0}) + \tfrac12\,\w^T\w \ \ge 0,
\end{equation}
which vanishes only for $q_{e,0}=1,\ \w=\vect 0$. From \eqref{eq:qdot},
$\dot q_{e,0} = -\tfrac12\q_{e,v}^T\w$, hence
\begin{equation}
  \dot V = K_p\,\q_{e,v}^T\w + \w^T\left(-K_p\q_{e,v} - K_d\w\right) = -K_d\norm{\w}^2 \le 0 .
\end{equation}
By LaSalle's invariance principle, the only invariant set with
$\dot V = 0$ is $\w=\vect 0,\ \q_{e,v}=\vect 0$: the equilibrium is asymptotically
stable. With constant disturbances and no integral action a small
steady-state error remains ($\approx 0.003^\circ$ in the simulator).

\begin{incode}
\code{controller.QuaternionPDController.compute(q, omega, h\_rw, q\_target)}.
\end{incode}

\begin{improvement}
\begin{itemize}[nosep]
  \item Integral term (PID) to remove the steady-state error.
  \item Limiting the manoeuvre rate (\emph{rate limiting}) for
        large angles, or \emph{bang-coast-bang} profiles.
  \item Optimal control (LQR/MPC) as a further baseline for RL.
  \item Momentum management: combined wheel + magnetorquer control.
\end{itemize}
\end{improvement}

% =============================================================================
\section{Model verification}
% =============================================================================

The file \code{tests/test\_physics.py} checks the fundamental physical properties
(\code{python -m pytest -q}):
\begin{center}
\begin{tabular}{p{5.2cm}p{9cm}}
\toprule
Test & Property checked\\
\midrule
\code{inertia\_matrix\_3U} & $J$ matches the analytical formula and is symmetric.\\
\code{torque\_free\_conservation} & Free motion for 100 s: constant $\vect H_I$ ($<10^{-9}$), constant $K$ ($<10^{-7}$ relative), $\norm{\q}=1$.\\
\code{internal\_torques\_conserve...} & Random wheel torques: total $\vect H_I$ conserved (Prop.~\ref{prop:H}).\\
\code{wheel\_saturation\_is\_exact...} & At maximum torque the wheels reach exactly 6000 RPM without exceeding it; $\vect H_I$ conserved.\\
\code{power\_and\_energy\_idle} & At zero torque $P=N P_s$ and $E = N P_s t$.\\
\code{quaternion\_kinematics\_analytic} & Uniform rotation about a principal axis: $\q(t) = [\cos\frac{\omega t}{2},0,0,\sin\frac{\omega t}{2}]$.\\
\code{gravity\_gradient\_zero...} & $\T_{gg}=\vect 0$ with nadir along a principal axis.\\
\code{pd\_controller\_converges} & From $90^\circ$ with disturbances: error $<0.1^\circ$ after 80 s.\\
\code{torque\_rate\_limit} & Step command: the applied torque changes by at most $\dot T_{max}\Delta t$ per step.\\
\bottomrule
\end{tabular}
\end{center}

\begin{improvement}
\begin{itemize}[nosep]
  \item Convergence study in $\Delta t$ (check the 4th order of RK4).
  \item Comparison with a reference simulator (e.g. Basilisk, 42) on
        identical scenarios.
  \item Tests of the intermediate-axis instability and of free precession
        against the analytical Euler--Poinsot solutions.
\end{itemize}
\end{improvement}

% =============================================================================
\section{Default parameters}
% =============================================================================
\begin{center}
\begin{tabular}{lll}
\toprule
Parameter & Value & Notes\\
\midrule
Mass, dimensions & 3 kg, $10\times10\times30$ cm & 3U CubeSat\\
Products of inertia & $(2,-1,1.5)\cdot10^{-5}$ kg\,m$^2$ & $(J_{xy},J_{xz},J_{yz})$\\
Centre-of-mass offset & $[2,-1,10]$ mm & SRP lever arm\\
$I_{rw}$ & $1.5\cdot10^{-5}$ kg\,m$^2$ & CubeSat-class wheel\\
$\Omega_{max}$, $T_{max}$ & 6000 RPM, 2 mN\,m & \\
Torque rate $\dot T_{max}$ & 8 mN\,m/s & motor driver limit\\
$k_1$, $k_2$, $P_s$ & 50 W/(N\,m), 1.2, 0.15 W & power model\\
Orbit & 500 km, $i=51.6^\circ$ & circular\\
Residual dipole & $[5,-3,10]$ mA\,m$^2$ & \\
$\rho_s$, $\rho_d$ & 0.1, 0.3 & \\
$\Delta t$ & 0.05 s & RK4\\
PD & $\omega_n=0.4$ rad/s, $\zeta=0.9$ & \\
\bottomrule
\end{tabular}
\end{center}
All parameters are defined as a \code{dataclass} in \code{satsim/config.py}.

% =============================================================================
\section{Towards Reinforcement Learning}
% =============================================================================

The engine is already structured as a discrete-time environment:
$\vect x_{n+1} = \Phi_{RK4}(\vect x_n, \vect u_n)$. A possible formulation
of the MDP problem (the formulation actually adopted, version by version, is
documented in \code{rl/README.md}):
\begin{itemize}
  \item \textbf{Action}: $\vect a\in[-1,1]^N$, with $\vect T_{rw} = T_{max}\,\vect a$.
  \item \textbf{Observation}: $[\q_e,\ \w,\ \vect\Omega/\Omega_{max}]$
        (possibly with sensor noise). It is convenient to force $q_{e,0}\ge0$
        to remove the sign ambiguity.
  \item \textbf{Reward}, for example:
        \begin{equation}
          r_n = -\alpha\,\theta_{err} - \beta\norm{\w}^2 - \gamma\,P_{tot}\Delta t
                - \delta\sum_i \max\!\left(0, \tfrac{|\Omega_i|}{\Omega_{max}} - 0.8\right)^2
                + r_{bonus}\,\mathbb{1}[\theta_{err} < \theta_{tol}],
        \end{equation}
        which balances precision, stability, energy and saturation margin.
  \item \textbf{Episodes}: random initial conditions (attitude uniform on
        $SO(3)$, \code{random\_quaternion}), random tumbling rates,
        randomised parameters ($J$, $I_{rw}$, residual dipole).
  \item \textbf{Baseline}: the PD of \S\ref{sec:pd}, to be compared on settling
        time, energy consumed and steady-state error.
\end{itemize}

% =============================================================================
\section*{References}
\addcontentsline{toc}{section}{References}
\begin{enumerate}[label={[\arabic*]}]
  \item F.~L.~Markley, J.~L.~Crassidis, \emph{Fundamentals of Spacecraft
        Attitude Determination and Control}, Springer, 2014.
  \item J.~R.~Wertz (ed.), \emph{Spacecraft Attitude Determination and
        Control}, Kluwer, 1978.
  \item M.~J.~Sidi, \emph{Spacecraft Dynamics and Control: A Practical
        Engineering Approach}, Cambridge University Press, 1997.
  \item B.~Wie, \emph{Space Vehicle Dynamics and Control}, 2nd ed., AIAA, 2008.
  \item H.~Schaub, J.~L.~Junkins, \emph{Analytical Mechanics of Space Systems},
        4th ed., AIAA, 2018.
  \item J.~Sol\`a, \emph{Quaternion kinematics for the error-state Kalman
        filter}, arXiv:1711.02508, 2017.
  \item O.~Montenbruck, E.~Gill, \emph{Satellite Orbits: Models, Methods and
        Applications}, Springer, 2000.
\end{enumerate}

\end{document}
